\section{Four-vertex complement lifts}\label{sec:complements}

\begin{theorem}[D4: complement-path lift]\label{thm:complementlift}
Let a simple graph $H$ of order $n\ge3$ have a given legal doubled
colouring and $\Delta(H)\le n-3$. Let $uv$ be a nonedge whose
endpoints have different ordered colour-degree pairs. In the
complement replace $uv$ by the five-edge path $uabcdv$ and take
the complement again, obtaining $H^+$. Its colouring can be extended
preserving all old edge weights and increasing every old red and
blue degree by four. Moreover $\Delta(H^+)\le |H^+|-3$, and
all five edges of the new complement path have endpoints with
different ordered colour-degree pairs.
\end{theorem}
\begin{proof}
First suppose $B_H(u)\ne B_H(v)$. Give the new edge $uv$ weight
zero. For every old vertex $w$ give $wb,wc$ weight 2 and $wa,wd$
weight zero whenever those edges exist; $ua$ and $vd$ are absent.
Among new vertices the actual edges are $ac,bd,ad$; give them
weights $1,2,0$. Old vertices different from $u,v$ gain $(4,4)$
from their four new edges. At $u,v$, the extra old--old edge $uv$
replaces the missing blue contribution at $ua,vd$, so they also
gain $(4,4)$. All old comparisons persist and $uv$ is legal.

The new counts are
\[
 a:(1,2n+1),\quad b:(2n+2,0),\quad
 c:(2n+1,1),\quad d:(2,2n).
\]
Every old final colour count is at most $2(n-3)+4=2n-2$.
Consequently pure red old--$b,c$ edges and pure blue old--$a,d$
edges are legal. The three actual new--new edges are legal by their
displayed counts. The five missing path edges have different ordered
pairs as well, by the same strict bounds and the distinct new pairs.
Each old degree increases by four, and every new degree is $n+1$,
which proves the maximum-degree assertion. If only the red counts
of $u,v$ differ, complement the old colours temporarily, apply this
construction, and complement the entire result back. Old weights
are again preserved and all stated properties remain true.
\end{proof}

Different boundary pairs are a real hypothesis. On $C_5$ the cyclic
word $(0,0,1,1,2)$ is legal, but two nonadjacent vertices have the
same pair $(2,2)$. A construction adding the same $(4,4)$ to each
and adding an actual edge between them cannot make that edge legal
with any weight. This restricts this lift interface, not the graph.

\subsection{An explicit nonregular infinite family}

For odd $k\ge3$ let
$F_k=C_{2k+1}+02$, using cyclic labels $0,\ldots,2k$, and
$G_k=\overline{F_k}$. Its two chord endpoints have degree $2k-3$
and the other vertices degree $2k-2$; its size is $2k^2-k-2$.
For $k=3$, in the edge order
\[
 03,04,05,13,14,15,16,24,25,26,35,36,46
\]
use weights
\[
 2,2,2,0,0,0,0,0,2,0,1,0,1.
\]
The vertex red degrees are $(6,0,2,3,3,5,1)$ and blue degrees
$(0,8,4,5,5,3,7)$, verifying actual legality. The complement edge
$60$ has different blue endpoint counts. Replacing it by a
five-edge path produces $G_5$ with new labels extending the cyclic
order, keeping chord $02$. Repeat on the last new edge incident
with 0. Its new endpoint blue degree is $2n$, greater than the old
endpoint final bound $2n-2$, so the interface stays available.
The lift proves colourability of every $G_k$, $k$ odd, by induction.

This family also limits a parity-refined defective partition. Suppose
one requires a partition $(A,B)$ with
\[
 \Delta(G_k[A])\le k-1,\quad \Delta(G_k[B])\le k-2,
 \quad R|_A\text{ even},\quad R|_B\text{ odd}.
\]
The degree conditions imply minimum induced complement degrees
$|A|-k$ on $A$ and $|B|-k+1$ on $B$.
Only two vertices of $F_k$ have degree three; hence these bounds
force $|A|\le k+2$, $|B|\le k+1$, and
$|B|\in\{k-1,k,k+1\}$. A minimum-degree-two induced subset of
$F_k$ with fewer than $2k$ vertices can only be its triangle $012$:
if it contains any vertex on the other arc from 2 to 0, its two
neighbours are forced, propagating along that entire arc of $2k$
vertices. If $|B|=k-1$, then $A$ has size $k+2\ge5$ and induced
minimum degree two, impossible; if $|B|=k+1$, the same holds for
$B$ of size at least four. Therefore $|B|=k$ is odd, contradicting
the even red handshake sum. The degree partition itself exists:
$B=\{0,\ldots,k-1\}$ and its complement both have no isolates in
$F_k$, which verifies the required complement degree bounds. It is
the appended parity prescription that fails. The legal colourings
above show that this is not a conjecture counterexample.

\subsection{Lifting complements of cubic graphs}

Let $F\ne K_4$ be a connected edge-transitive cubic simple graph
of order at least six and $H=\overline F$. Each component of $H$
is regular of degree $n-4\ge2$, so it has a legal doubled colouring
by~\cite{GOW26}. Some edge of $F$ has distinct endpoint pairs:
otherwise connectedness of $F$ would make the pair constant on all
vertices, contradicting legality of any actual edge of $H$.
An automorphism carries that complement interface to any selected
edge of $F$. Repeated use of Theorem~\ref{thm:complementlift}
therefore handles replacing that selected edge of $F$ by a path
of length $4t+1$ and then complementing, for any $t\ge0$. At
$t>0$ these graphs are nonregular. Edge-transitivity is used to
control the selected initial edge and is not silently omitted.

\subsection{Restrictions in complement degree two or three}

Suppose $G$ is a vertex-minimal counterexample to the general
conjecture and every vertex of $F=\overline G$ has degree two or
three. If $n\ge11$, $F$ cannot contain a path $uabcdv$ whose four
internal vertices have degree two, whose endpoint degrees differ,
and with $uv\notin E(F)$. Delete the four internal vertices and
replace that path by $uv$. The resulting complement $H$ has order
$N=n-4\ge7$, degree between $N-4$ and $N-3$, and is connected
because $2(N-3)>N$. Minimality gives a legal colouring of $H$.
The two endpoint ordinary degrees differ, so their ordered colour
pairs differ. The lift reconstructs $G$, a contradiction.
It follows that no path starting at an $F$-degree-three vertex can
continue through six consecutive distinct degree-two vertices:
use its first five-edge subpath, whose endpoint degrees are three
and two; the sixth vertex ensures the missing closing edge.

There is also no cycle component of $F$ of length at least nine.
Shorten that cycle by four. Its length remains at least five, so
its complement is connected (and is completely joined to any other
complement components). Minimality supplies a colouring of the
smaller actual graph. Some edge of the shortened complement cycle
has distinct endpoint pairs; otherwise every vertex of that cycle
would have the same pair, contradicting an actual edge within its
complement. Rotate that cycle component to place the interface at
the edge to be lengthened and apply the lift. Finally
Theorem~\ref{thm:dense} covers this entire complement-degree class
when $n\ge38$ since $q\le3$. Any hypothetical counterexample in
this class therefore has $n<38$. No exhaustion of that residual
finite range, or reduction of all graphs to this class, is claimed.
